\hypertarget{appendix-a427ce8ce8a6}{}
\section*{Bitwise OR}
\section{Bitwise OR}\label{bitwise-or}

\textbf{Curriculum depth:} Introduced · \textbf{First seen:} C001 · \textbf{Latest development:} C001

\subsection{Definition and mechanism}\label{definition-and-mechanism}

Sets selected lanes while preserving lanes paired with 0; it cannot clear an existing 1.

Per lane: only 0\textbar0 is 0; the other three cases are 1. \texttt{0xAD\ \textbar{}\ 0x40\ =\ 0xED} sets locked. Repeating the operation leaves the same result: OR is idempotent. It cannot clear an existing 1, so replacing a field requires a clear step.

\subsection{Prerequisites}\label{prerequisites}

\begin{itemize}
\tightlist
\item
  \hyperref[appendix-24475e1596a8]{Boolean algebra} --- The logical domain is false/true.
\end{itemize}

\subsection{Used by}\label{used-by}

\begin{itemize}
\tightlist
\item
  \hyperref[appendix-1f9963e1a5d1]{Bit packing} --- uses bitwise or as part of its representation or reasoning.
\item
  \hyperref[appendix-fe6b457d68d8]{Logical versus bitwise operators} --- uses bitwise or as part of its representation or reasoning.
\item
  \hyperref[appendix-b5c46b2d94ff]{Carry-lookahead and parallel-prefix adders} --- uses bitwise or as part of its representation or reasoning.
\end{itemize}

\subsection{Recall and next depth}\label{recall-and-next-depth}

Explain bitwise or without the label, work the example above, then state the assumption whose removal would change the result.

\textbf{Pending:} flag sets and concurrent updates.

\subsection{Source and navigation}\label{source-and-navigation}

\href{Continuous\%20Notes/Computing\%20Foundations\%20-\%20From\%20Bits\%20to\%20C++\%20and\%20Systems\%20APIs\#From\%20a\%20signal\%20to\%20an\%20interpretation}{C001 teaching chapter} · \hyperref[appendix-fd090d982a15]{Part definitions and mechanisms} · \href{}{Dependency Map (02-Foundations-II, appendix-4ad47bdd269a)}
